\section{Detailed scope and limitations}
\label{app:detailed-limitations}

The conclusions are bounded by six aspects of the evaluation: the compact-evidence task regime, the fitted sparse interfaces, the candidate families used for path identification, the interface-specific scope of the full-hidden confirmation and matched follow-up, answer-conditioned native-channel selection, and the annotation workflow. These boundaries restrict how the results should be generalized without changing the central finding within the evaluated setting.

\textbf{Evaluation regime.} The study targets compact-evidence visual question answering examples in which an answer-bearing image region can be localized and corrupted. It does not estimate how often such evidence occurs across the full OK-VQA distribution, nor does it extend the same conclusions to diffuse scene reasoning. The 61-example expanded and 52-example independently selected sets support the regional and hidden-position comparisons. The stronger content, direction, location, sparse-tool, and native-channel tests use a separate 60-image development and 240-image reused confirmation split; this split is confirmatory with respect to frozen choices but is not previously unseen project data.

\textbf{Sparse tools remain model-specific approximations.} The per-layer and cross-layer transcoders differ in training, target layers, and reconstruction quality. Their evidence-induced change-reconstruction errors exceed one in the tested answer positions, so the absence of a positive sparse-component answer effect does not establish that the base models lack sparse computation. Conversely, the observed structural differences describe the fitted interfaces and cannot be attributed uniquely to the base architectures.

\textbf{The tested compact paths are not exhaustive.} Native-channel tests cover two frozen answer-conditioned selection rules, while the shared-subspace tests cover uncentered linear projections of ranks 1--128 at the initial model-selected and later matched late-answer interfaces. Each basis is fitted separately within a model and shared only across that model's examples. Failure of these candidates excludes the corresponding channel groups and example-shared linear subspaces; it does not exclude higher-rank, nonlinear, layer-spanning, or sample-specific objects. Principal components are fitted from clean-minus-masked differences and are analysis objects rather than variables explicitly represented by the base models.

\textbf{The full-hidden finding belongs to the evaluated interfaces.} On 60 reviewed examples, restoring the complete hidden state at LLaVA layer 30 and four answer-adjacent positions recovers most of the clean-minus-masked sequence-score difference, while replacing its clean value with the masked value removes most of that difference. The later matched follow-up tests Gemma layer 32, Qwen layer 26, and LLaVA layer 30 at the same relative network depth and position type. Only the tested LLaVA interface passes all four bounds, but the rule was chosen after the LLaVA result and the 60 examples were not jointly unseen at that point. Different architectures, tokenization, and possible alternative layers remain. The finding does not establish minimum dimension, necessity of every coordinate, or absence of another successful interface in Gemma or Qwen. The confirmation masks were produced by GPT-assisted pre-annotation and reviewed individually without model results; 60 were accepted without correction and one was rejected before evaluation.

\textbf{Native-channel selection is answer-conditioned.} The native-channel experiment freezes the ranking rule and retained fraction on development images, but applies that rule to each confirmation image using its known correct answer. Both the positive-only links and the Qwen positive-and-negative link are therefore diagnostics of answer-conditioned, image-specific support sets. Establishing an answer-independent circuit would require selecting channels without using the evaluated image's target answer or confirming a shared channel set on separate images.

\textbf{Annotation and wording.} Evidence masks come from a single-pass annotation workflow, so independent inter-annotator agreement is unavailable. Same-area random, shifted, shuffled, dilation, erosion, blur, inpainting, and patch-shuffle controls reduce explanations based on one exact polygon or corruption operator but do not eliminate all mask uncertainty. Semantically similar prompts can also change which sparse identities or hidden locations appear most salient. We therefore claim robustness of the localized causal question, not invariance of one internal coordinate across wording.

The complete intervention tables, candidate-selection rules, sparse-tool settings, shared-subspace sweep, and additional mask and wording controls are given in Appendices~\ref{app:detailed-results}, \ref{app:sparse-lens-details}, \ref{app:shared-subspace-details}, \ref{app:regional-controls}, and~\ref{app:prompt-text}. Within the stated evaluation regime, the three models show an evidence-specific answer-score response and the frozen LLaVA full-hidden interface passes independent bidirectional confirmation.
